\subsection{Ước lượng}

\subsubsection{Khoảng tin cậy cho giá trị trung bình}

\begin{exmp}
  Nếu $X_1$, $X_2$, $\ldots\,$ là các biến ngẫu nhiên i.i.d Gaussian hoặc bị
  chặn trong $[-1, 1]$ thì \begin{equation}
    \frac{ \sum_{i = 1}^n X_i }{ n } \pm 1.71 \sqrt{ \frac{ \log\log(2n) + 0.72
    \log ( 10.4 / \alpha ) }{ n } }
  \end{equation} là một $(1 - \alpha)$ confidence sequence cho $\mu$, tương tự
  như \begin{equation}
    \bigcap_{s \le n} \frac{ \sum_{i = 1}^n X_i }{ n } \pm 1.71 \sqrt{ \frac{
      \log\log(2n) + 0.72 \log ( 10.4 / \alpha ) }{ n } }.
  \end{equation}
  Ta thấy rằng, độ rộng của $(1 - \alpha)$ confidence sequence lớn hơn so với
  khoảng tin cậy tương ứng nhưng chênh lệch không thực sự ảnh hưởng nhiều. Để
  thấy sự khác biệt giữa hai kết quả này. Ta dùng đến tỉ lệ bao phủ sai tích
  lũy (cummulative miscoverage rate). Tỉ lệ này ứng với mỗi $t$ chính là tỉ lệ
  tồn $t$ khoảng tin cậy đầu tiên bị bao phủ sai với các lần chạy. Do khoảng
  tin cậy chỉ đảm bảo sai lầm ở các lần chạy ở một thời điểm $t$ là dưới
  $\alpha$ nên khi xét đồng thời $t$ khoảng tin cậy thì tỉ lệ bao phủ đồng thời
  của $t$ khoảng tin cậy gần như bằng $0$ (hay tỉ lệ bao phủ sai gần như là
  $1$). Tuy nhiên, theo định nghĩa của confidence sequence, tỉ lệ bao phủ sai
  tích lũy luôn được đảm bảo dưới $\alpha$ do tính đồng thời của nó. Điều này
  cũng cho thấy tính hợp lệ tại mọi thời điểm của confidence sequence so với
  khoảng tin cậy. (Xem minh họa tại trang số 8, nửa bài giảng sau.)
\end{exmp}

\subsubsection{Khoảng tin cậy cho phân vị}

\paragraph{Một phân vị cố định}

\begin{exmp}
  Đặt \begin{equation*}
    u_t := \sqrt{ \frac{ 0.73 \log\log(2.04t) + 0.52 \log(9.97/\alpha) } { t }
    }
  \end{equation*} Khi đó, \begin{equation*}
    \prob(\forall t \in \mathbb{N} : \hat{Q}_t(1/ 2 - u_t) \le Q(1/2) \le
    \hat{Q}_t(1/2 + u_t)) \ge 1 - \alpha.
  \end{equation*} Tức là \begin{equation*}
    [\hat{Q}_t(1/ 2 - u_t), \hat{Q}_t(1/2 + u_t))], \quad \forall t
  \end{equation*} là một $(1 - \alpha)$ confidence sequence cho phân vị $1/2$
(trung vị).

\end{exmp}


\paragraph{Toàn bộ phân vị đồng thời}

\begin{exmp}
  Đặt \begin{equation*}
    u_t := \sqrt{ \frac{ \log\log(et) + 0.75 \log(34 / \alpha) } { t } }.
  \end{equation*} Khi đó, \begin{equation*}
    \prob(\forall t \in \mathbb{N}, p \in (0, 1): \hat{Q}_t(p - u_t) \le Q(p)
    \le \hat{Q}_t(p + u_t) ) \ge 1 - \alpha.
  \end{equation*} Tức là ta có đồng thời các khoảng tin cậy cho từng thời điểm
  với cùng độ rộng cho tất cả phân vị khác nhau. Và chúng đều tạo thành một $(1
  - \alpha)$ confidence sequence đồng thời cho tất cả các phân vị.
\end{exmp}

\subsubsection{Ước lượng ma trận hiệp phương sai tuần tự}

Xét biến ngẫu nhiên $X \in \mathbb{R}^d$ với kỳ vọng $\expect[X] = 0$ và điều
kiện bị chặn $|X_i| \le b$. Sai số hội tụ của ước lượng $\hat{\Sigma}_n$ so với
ma trận thực $\Sigma$ được chặn trên với xác suất cao như sau:
\begin{equation}
  \big\| \hat{\Sigma}_n - \Sigma \big\|_{\text{op}} \lesssim \sqrt{\frac{b^2
  \log(d \log n)}{n}} + \frac{b^2 \log(d \log n)}{n}.
\end{equation}

\subsubsection{Ước lượng qua cá cược}

Cho $X_1$, $X_2$, $\ldots\,$ là các biến ngẫu nhiên i.i.d thuộc $[0, 1]$ với
trung bình $\mu$. Ta thiết kế trò chơi như
\algorithmname~\ref{alg:esitmate-game}. Khi đó đặt \begin{equation*}
  C_t := \{ m \in [0, 1]: K^m_t < 1/\alpha \}
\end{equation*} với ý nghĩa rằng nhà khoa học không thắng được nhiều tiền. Ta
có định lý (có thể chứng minh thông qua kiểm định qua cá cược) nói rằng
$(C_t)_{t \ge 1}$ là một confidence sequence cho $\mu$ với mọi chiến lược đặt
cược.

\begin{algorithm}[htbp]
  \caption{Ước lượng qua cá cược}
  \label{alg:esitmate-game}

  Khởi tạo vốn ban đầu $K^m_0 = 1$ cho mỗi trò chơi $m \in [0, 1]$\;

  \For{$t = 1, 2, \dots$}{
    Với mỗi $m \in [0, 1]$, nhà thống kê đặt cược $\lambda_t^m \in [ -1/( 1- m
    ), 1/m ]$\;
    Quan sát $X_t$\;
    Cập nhật tài sản trong game $m$ thành $K^m_t = K^m_{t - 1} \cdot (1 +
    \lambda^m_t(X_t - m))$\;
  }
\end{algorithm}

Khi đó, vấn đề trở thành ta cần một chiến lược đặt cược để các tập hợp này càng
nhỏ càng tốt. Một trong những chiến lược là GRAPA (Grow Rate Adaptive to the
Particular Alternative), ta đặt cược như sau: \begin{equation}
  \lambda^m_t(P) := \arg\max_{\lambda \in [-1/(1 - m), 1/m]} \expect_P[\log(1 +
  \lambda(X_t - m)) \mid \mathcal{F}_{t - 1}].
\end{equation} Nhưng ta không biết $P$, lời giải xấp xỉ: lấy đạo hàm theo biến
$\lambda$, cho đạo hàm bằng 0 (áp dụng điều kiện KKT), khai triển Taylor và
thay thế bằng các ước lượng thực nghiệm. Khi đó ta có giá trị cho khoản cược là
\begin{equation*}
  \lambda^m_t = \frac{\hat\mu_t - m}{ \hat\sigma_t^2 + (\hat\mu_t - m)^2 },
\end{equation*} trong đó các $\mu_t$ và $\sigma_t^2$ được tính từ $t - 1$ mẫu
đầu.

\subsection*{Chiến lược cá cược 2: Phương pháp hỗn hợp (Mixture)}

Chiến lược thứ hai là Mixture. Chiến lược này được củng cố và phát triển thông
qua hàng loạt các nghiên cứu gần đây, tiêu biểu như Waudby-Smith \& Ramdas
(2024), Jun \& Orabona (2024), và Shekhar \& Ramdas (2024). Theo đó, confidence
sequence được định nghĩa dựa trên một tích phân hỗn hợp như sau:
\begin{equation}
  C_t := \Biggl\{ m \in [0,1] : \int_{-1/(1-m)}^{1/m} \prod_{i=1}^t \big(1 +
  \lambda(X_i - m)\big) \, d\nu(\lambda) < \frac{1}{\alpha} \Biggr\}
\end{equation} Về mặt triển khai thực tiễn, việc tính toán tích phân liên tục
thường được thay thế bằng phương pháp rời rạc hóa (discretize) không gian của
hỗn hợp. Khía cạnh đáng chú ý là phép xấp xỉ này hoàn toàn không làm ảnh hưởng
đến tính hợp lệ thống kê (validity) của confidence sequence.

Tuy nhiên, nhằm bảo toàn sức mạnh kiểm định (power) khi lượng dữ liệu tăng lên,
lưới rời rạc hóa của hỗn hợp bắt buộc phải được tinh chỉnh mịn hơn (finer) theo
thời gian. Cuối cùng, để đánh giá và phân tích hiệu năng của chiến lược này, ta
có thể thiết lập mối liên hệ trực tiếp với bài toán Tối ưu hóa Danh mục Đầu tư
của Cover (Cover's Portfolio Optimization). Dựa trên cơ sở lý thuyết đó, việc
sử dụng một phân phối hỗn hợp đều (uniform mixture) thường là phương án đủ để
đảm bảo tính tối ưu.

\subsubsection{}

Trong suy luận nhân quả và phân tích thực nghiệm, sự kết hợp giữa Tác động can
thiệp trung bình (ATE) và asymptotic confidence sequence (AsympCS) mang lại một
giải pháp đột phá cho bài toán đánh giá dữ liệu liên tục.

Trong khi ATE đóng vai trò là thước đo cốt lõi giúp xác định hiệu quả trung
bình thực sự của một sự can thiệp (như một tính năng công nghệ mới hay một phác
đồ điều trị) trên toàn bộ quần thể, thì việc liên tục theo dõi ATE để ra quyết
định sớm lại rất dễ vi phạm các nguyên tắc thống kê truyền thống. Đây là lúc
AsympCS trở thành mảnh ghép hoàn hảo. Bằng cách nới lỏng các ràng buộc toán học
khắt khe tuyệt đối thành dạng xấp xỉ tiệm cận (thường dựa trên Định lý giới hạn
trung tâm), AsympCS cung cấp một dải an toàn cho phép bạn liên tục "nhìn" vào
kết quả ATE theo thời gian thực (anytime-valid inference). Nhờ đó, các nhà
nghiên cứu và kỹ sư dữ liệu có thể linh hoạt dừng thử nghiệm sớm ngay khi hiệu
quả của can thiệp đã đủ rõ ràng, vừa tiết kiệm thời gian, chi phí, lại vừa đảm
bảo tính hợp lệ thống kê chặt chẽ trong dài hạn.
